\documentclass[11pt,a4paper]{article}
\usepackage{amsmath,amssymb,amsthm}
\usepackage{mathtools}
\usepackage[margin=2.3cm]{geometry}
\usepackage{microtype}
\usepackage{booktabs}
\usepackage{array}

\theoremstyle{plain}
\newtheorem{lemma}{Lemma}
\newtheorem{proposition}[lemma]{Proposition}
\newtheorem{corollary}[lemma]{Corollary}
\theoremstyle{definition}
\newtheorem{definition}[lemma]{Definition}
\newtheorem{remark}[lemma]{Remark}
\newtheorem{question}[lemma]{Question}

\newcommand{\Zp}{\mathbb{Z}}
\newcommand{\R}{R}
\newcommand{\rowsp}{\operatorname{rowsp}}
\newcommand{\colsp}{\operatorname{colsp}}
\newcommand{\im}{\operatorname{im}}
\newcommand{\rank}{\operatorname{rank}}
\newcommand{\sat}[1]{\overline{#1}}
\newcommand{\T}{^{\mathsf T}}
\newcommand{\Fp}{\mathbb{F}_p}

\title{Scratch notes: the chain-ring factorisation, and\\
the class as a product of constructions}
\author{(working notes for the bounded-rank-width paper)}
\date{\today}

\begin{document}
\maketitle

\noindent\textbf{Purpose.} These notes settle whether the tree-merge
machinery of the paper survives the passage from the field $\Fp$ to the
ring $\R = \Zp/p^{d}\Zp$ (the extension ``Idea~2'': class-$2$ groups whose
\emph{center} has exponent $p^{d}$, so the cocycle becomes an
$\R$-bilinear form). The crux is the two-sided factorisation
lemma. \textbf{Headline: the lemma is false over $\R$ as stated
(counterexample at $r=1$), but becomes true again once the carried
interface is a \emph{free} basis of the saturation of the row/column
space. The valuation that broke the naive version is absorbed into the
factor, not the interface.} All claims below are checked by exhaustive
machine search over $\R\in\{\Zp/4,\Zp/8,\Zp/9\}$ (see
\S\ref{sec:machine}).

\section{Setup}

Throughout $\R=\Zp/p^{d}\Zp$ is a finite chain ring: a local principal
ideal ring whose ideals are the chain $\R\supsetneq p\R\supsetneq\dots
\supsetneq p^{d}\R=0$. Every element is $u\,p^{s}$ with $u$ a unit and
$0\le s\le d$ its \emph{valuation} $v(\cdot)$; the units are exactly the
valuation-$0$ elements. Modules are finitely generated. A submodule
$M\subseteq \R^{m}$ is \emph{pure} (a \emph{direct summand}, since $\R$ is
quasi-Frobenius and $\R^m$ is free) if $\R^{m}=M\oplus M'$ for some $M'$;
the \emph{saturation} $\sat M$ is the smallest pure submodule containing
$M$. We write $\rowsp(V)$ for the row module of a matrix $V$ (the
$\R$-span of its rows) and $\colsp(V)$ for its column module.

Recall the field statement that drives the tree construction.

\begin{lemma}[two-sided factorisation over a field]\label{lem:field}
Let $X\in\Fp^{m\times m'}$, $V\in\Fp^{r\times m'}$, $W\in\Fp^{r\times m}$.
If every row of $X$ lies in $\rowsp(V)$ and every column of $X$ lies in
$\colsp(W\T)$, then $X=W\T Q\,V$ for some $Q\in\Fp^{r\times r}$.
\end{lemma}

\begin{proof}
Rows of $X$ in $\rowsp(V)$ give $X=AV$ for some $A$. Discarding dependent
rows of $V$ we may assume its rows are independent, so $V$ has a right
inverse $Y$ ($VY=I_r$); then $A=XY$, so every column of $A$ is a
combination of columns of $X$, hence lies in $\colsp(W\T)$, and
$A=W\T Q$. \qed
\end{proof}

The single ring-sensitive step is ``$V$ has a right inverse''. Over $\R$ a
matrix with independent (even minimal) rows need not be right-invertible:
$V=(p)$ over $\Zp/p^{2}$ has $\rowsp(V)=p\R\cong\Zp/p$, which is not a
direct summand, and $V$ has no right inverse.

\section{Failure over the chain ring}

\begin{proposition}[the lemma is false over $\R$, already at $r=1$]
\label{prop:fail}
Over $\R=\Zp/4$ put
\[
  V=W=\begin{pmatrix}2&0\end{pmatrix},\qquad
  X=\begin{pmatrix}2&0\\0&0\end{pmatrix}.
\]
Then every row of $X$ lies in $\rowsp(V)=\{(0,0),(2,0)\}$ and every column
of $X$ lies in $\colsp(W\T)=\{(0,0)\T,(2,0)\T\}$, yet there is no scalar
$q\in\R$ with $X=W\T q\,V$.
\end{proposition}

\begin{proof}
The hypotheses are immediate. For the conclusion,
\[
  W\T q\,V
  = q\begin{pmatrix}2\\0\end{pmatrix}\begin{pmatrix}2&0\end{pmatrix}
  = q\begin{pmatrix}4&0\\0&0\end{pmatrix}
  = \begin{pmatrix}0&0\\0&0\end{pmatrix}\neq X
  \qquad\text{for every }q,
\]
because $4=0$ in $\R$. \qed
\end{proof}

\begin{remark}[what goes wrong: valuation is lost in the interface]
The interface module $\rowsp(V)=p\R$ is generated by the
\emph{valuation-$1$} vector $(2,0)=p\cdot(1,0)$. A product $W\T q V$ pairs
two valuation-$1$ rows and therefore only ever produces entries of
valuation $\ge 2$, i.e.\ $0$ modulo $p^{2}$. The target $X$ has a
valuation-$1$ entry, so it cannot be hit. Minimal generating sets
(``module rank'', the number of invariant factors) are blind to this: the
row module has rank $1$, but the naive rank-$1$ interface has already
spent its valuation budget. This is exactly the phenomenon that has no
analogue over a field, where every nonzero element is a unit.
\end{remark}

\section{The fix: saturate the interface}

The remedy is to carry not a minimal generating set of the row module but
a \emph{free basis of its saturation}, letting the valuation live in the
factor $Q$. Concretely, in Proposition~\ref{prop:fail} replace $V=(2,0)$
by the free basis $\sat V=(1,0)$ of $\sat{\rowsp(V)}=\R\cdot(1,0)$ (a
direct summand), and likewise $\sat W=(1,0)$. Then
\[
  X=\begin{pmatrix}2&0\\0&0\end{pmatrix}
   =\begin{pmatrix}1\\0\end{pmatrix}\,(2)\,\begin{pmatrix}1&0\end{pmatrix}
   =\sat W\T\,(2)\,\sat V,
\]
and the valuation-$1$ factor $q=2$ now sits in the middle, where it is
harmless. The general statement:

\begin{lemma}[two-sided factorisation over $\R$, saturated form]
\label{lem:ring}
Let $X\in\R^{m\times m'}$, and let $V\in\R^{r\times m'}$, $W\in\R^{r\times
m}$ be \emph{right-invertible} (equivalently, their rows are a free basis
of a direct summand; equivalently, some $r\times r$ minor is a unit). If
every row of $X$ lies in $\rowsp(V)$ and every column of $X$ lies in
$\colsp(W\T)$, then $X=W\T Q\,V$ for some $Q\in\R^{r\times r}$.
\end{lemma}

\begin{proof}
Verbatim the proof of Lemma~\ref{lem:field}: rows of $X$ in $\rowsp(V)$
give $X=AV$; the right inverse $Y$ of $V$ exists \emph{by hypothesis}, so
$A=XY$ has its columns among $\R$-combinations of columns of $X$, hence in
$\colsp(W\T)$, so $A=W\T Q$. The only place the field proof used more than
ring axioms was the existence of $Y$, which is now assumed. \qed
\end{proof}

So the content is entirely in \emph{choosing} the interfaces to be free.
That this is always possible, at no cost in width, is the next point.

\begin{lemma}[saturation is free of bounded rank]\label{lem:sat}
Let $M\subseteq\R^{m'}$ be generated by $r$ elements. Then $\sat M$ is a
free direct summand of $\R^{m'}$ of rank
\[
  \rho(M):=\dim_{\Fp}\bigl(M/pM\bigr)\ \le\ r ,
\]
and $M=\sat M\cdot D$ for a diagonal $D=\operatorname{diag}(p^{e_1},\dots,
p^{e_{\rho}})$ with $0\le e_i<d$ in a suitable free basis of $\sat M$.
\end{lemma}

\begin{proof}
Over the chain ring $\R$ (a PIR) the submodule $M$ has a Smith normal
form: $M=\rowsp\bigl(\operatorname{diag}(p^{e_1},\dots)\,U\bigr)$ for a
unimodular $U$ and exponents $0\le e_i<d$ (entries with $e_i\ge d$ vanish
and are dropped). The rows of $U$ corresponding to the $e_i<d$ span the
free summand $\sat M$; there are $\dim_{\Fp}M/pM$ of them by Nakayama, and
this is at most the number of generators $r$. Scaling the $i$-th basis
row by $p^{e_i}$ recovers $M$, which is the displayed factorisation. \qed
\end{proof}

\begin{corollary}[the merge factorisation survives]\label{cor:merge}
Let $X$ be a block of the cocycle whose row module and column module are
generated by $\le r$ elements, with all rows of $X$ in the row module and
all columns in the column module. Carrying the free saturated bases $\sat
V$ (rank $\rho\le r$) and $\sat W$ (rank $\rho'\le r$), one has
$X=\sat W\T\,Q\,\sat V$ for some $Q\in\R^{\rho'\times\rho}$.
\end{corollary}

\begin{proof}
By Lemma~\ref{lem:sat} the rows of $X$, lying in the row module
$\subseteq\sat{\rowsp}$, lie in $\rowsp(\sat V)$, and dually for columns;
$\sat V,\sat W$ are right-invertible, so Lemma~\ref{lem:ring} applies. \qed
\end{proof}

\section{Consequence for the exponent-$p^{d}$ extension}

Write $A=G/Z$ for the abelian quotient and let the center $Z$ have
exponent $p^{d}$, so the commutator pairing is an $\R$-bilinear map
$C\colon (A\otimes\R)\times(A\otimes\R)\to\R^{k}$ and the ``form'' $B$ has
entries in $\R$. Define the \emph{module cut-rank} of a layout as the
maximum, over cuts $S$, of $\rho(\text{crossing block at }S)$ --- the
number of invariant factors, equivalently $\dim_{\Fp}(\text{row
module}\otimes\Fp)$ --- \emph{measured on saturations}, i.e.\ the free
rank of the saturated interface. Then:

\begin{itemize}
\item \textbf{One-sided streaming is unconditional.} The word/linear
  automaton (Theorem~1 of the paper) uses only that the restriction of the
  row module to earlier columns lies in the previous row module; this is a
  statement about module homomorphisms and holds over any ring. Carrying
  the saturated free basis $\sat V_t$ (rank $\le r$), with the update
  $\sat V_t=[\,T_t\,\sat V_{t-1}\mid v_t\,]$ solved over $\R$ and the
  valuations absorbed into $T_t$ and the read-off $R_t$, the exponent-$p^d$
  \emph{word} theorem lifts directly.
\item \textbf{The tree merge survives via Corollary~\ref{cor:merge}.} The
  two-sided factorisation --- the one place the field proof genuinely used
  a field --- holds provided the advice carries \emph{saturated} bases.
  This is the precise answer to ``does rank-width survive the passage from
  a field to $\Zp/p^{d}$'': yes, but the correct width is the free rank of
  the \emph{pure closure}, not the minimal number of generators, and the
  factor $Q$ (not the interface) carries the valuations.
\item \textbf{Everything stays bounded.} $\rho(M)\le r$
  (Lemma~\ref{lem:sat}), each carried register is an element of a free
  module of rank $\le r$ over $\R$, i.e.\ $\le rd$ digits mod $p$; the
  state count is $p^{O(rd)}$, matching the field case up to the factor $d$.
\end{itemize}

\begin{remark}[what still needs care in the full construction]
Corollary~\ref{cor:merge} fixes the algebraic crux. Two bookkeeping items
remain for the full write-up, neither ring-specific in spirit:
(i) a \emph{single} free interface of rank $r$ must serve every cut ---
saturations of successive row modules need not nest under column
restriction, so one carries a free module of the maximal rank and tracks
the (possibly non-unit) transition maps between cuts, which
Lemma~\ref{lem:sat} guarantees exist; (ii) the base-$p$ carries that
implement multiplication in $\R$ (a valuation-$1$ product feeding a higher
digit) are the same local, regular bookkeeping already used for the
abelian carries of ``Idea~1''. Both are routine given
Corollary~\ref{cor:merge}; the field$\to$ring obstruction was entirely in
the two-sided factorisation, and it is resolved.
\end{remark}

\section{Machine verification}\label{sec:machine}

All claims were checked by exhaustive/structured search over
$\R\in\{\Zp/4,\Zp/8,\Zp/9,\Zp/16\}$:
\begin{itemize}
\item \textbf{Naive lemma fails.} Searching arbitrary (structured $p$-scaled
  plus random) bases $V,W$ at $r=1,2$ finds counterexamples immediately;
  the minimal one over $\Zp/4$ is Proposition~\ref{prop:fail}.
\item \textbf{Saturated lemma holds.} Restricting $V,W$ to right-invertible
  bases, no counterexample occurs among $6.4\times10^{3}$ ($\Zp/4$, $r=1$),
  $4.1\times10^{5}$ ($\Zp/4$, $r=2$), $6.6\times10^{6}$ ($\Zp/8$, $r=2$)
  and $1.05\times10^{7}$ ($\Zp/9$, $r=2$) admissible blocks $X$ --- i.e.\
  Lemma~\ref{lem:ring} holds on every case tested.
\end{itemize}
Scripts: \texttt{scratchpad/twosided2.py} (finds the counterexamples),
\texttt{scratchpad/twosided3.py} (confirms the saturated form).

\section{The class as a product of constructions}
\label{sec:span}

With the chain-ring row secured, it is worth stepping back: the class is
not one construction but a \emph{product of independent axes}, each a
group-theoretic knob, and uniform automaticity is preserved for every
setting. Table~\ref{tab:axes} lists the axes; Table~\ref{tab:constructions}
places the named constructions as coordinates in that space.

\begin{table}[h]
\centering\small
\renewcommand{\arraystretch}{1.25}
\begin{tabular}{@{}p{2.0cm}p{2.5cm}p{2.4cm}p{4.4cm}p{2.7cm}@{}}
\toprule
Axis (knob) & Off & On & Group-theoretic meaning & Cost of ``on'' \\
\midrule
Layout & linear (word) & tree (branching) & path-like vs.\ laminar or
branching commutation & word $\to$ tree automaton \\
Width $r$ & $r=0$ (abelian) & $r\ge 1$ & commutator complexity across each
cut & $p^{\Theta(r)}$ states \\
Center & $k=0$ or fixed $k$ on a chain & distributed, growing & size and
placement of $G'$ & claim--verify, $p^{\Theta(r)}$ \\
Quotient exponent (Idea~1) & $\exp p$: $G/Z$ elem.\ abelian & arbitrary
abelian $G/Z$ & orders of the generators & base-$p$ carries (free) \\
Center exponent (Idea~2) & $\exp p$: field $\Fp$ & $\exp p^{d}$: ring
$\Zp/p^{d}$ & $G'$ elementary vs.\ higher & saturated width, $\times d$
digits \\
Composition & single form & $\bigoplus$ forms & central / direct products &
width $=\max$ of blocks \\
\bottomrule
\end{tabular}
\caption{The independent axes of the class. The whole product stays
uniformly automatic; the ``on'' end of each row is a strictly more general
group-theoretic construction, at the stated cost.}
\label{tab:axes}
\end{table}

The composition row is the crux of the model-theoretic point: central and
direct products are block-diagonal forms, and with a block-consecutive
layout a cut either splits within one block or crosses between blocks with
no commutators, so the width is the \emph{maximum} over blocks, not the
sum. One may therefore amalgamate arbitrarily many bounded-width building
blocks and never leave the class.

\begin{table}[h]
\centering\small
\renewcommand{\arraystretch}{1.2}
\begin{tabular}{@{}p{4.2cm}p{7.2cm}p{3.1cm}@{}}
\toprule
Construction & Setting & Status \\
\midrule
finite abelian $p$-group & $r=0$ & \checkmark\ base corner \\
extraspecial $p^{1+2m}$ & linear, matching, $k=1$, $\exp p$ &
\checkmark\ width $1$ \\
Heisenberg over $\Fp^{n}$ & fixed $k$, symplectic form & \checkmark \\
special $p$-groups & fixed $k$, form of width $\le r$ & \checkmark\ if
width bounded \\
$n$-qubit Pauli group & $p=2$, clique & \checkmark\ width $1$ \\
shared-center graph group & $[x_i,x_j]=y$ on edges, graph rank-width
$\le r$ & \checkmark \\
tree-indexed extraspecial & distributed center, laminar & \checkmark\ width
$1$ \\
point-target family & distributed, growing, non-laminar & \checkmark\ width
$1$ \\
central / direct products of the above & $\bigoplus$ forms & \checkmark\
width $=\max$ \\
higher-exponent quotient & Idea~1 on & \checkmark \\
higher-exponent commutator & Idea~2 on, over $\Zp/p^{d}$ & \checkmark\
saturated width \\
\midrule
ultraspecial (Heisenberg over $\mathbb F_{p^{m}}$) & field-multiplication
tensor & $\times$\ width $\Theta(m)$ \\
Mekler groups (private targets) & one free target per relation &
$\times$\ generically unbounded \\
relatively free class-$2$, $\exp p$ & fully free & $\times$\ non-automatic
(Nies--Stephan) \\
\bottomrule
\end{tabular}
\caption{Named constructions as settings of the axes. The three excluded
rows are exactly the model-theoretically wild objects; bounded rank-width
is the dividing line.}
\label{tab:constructions}
\end{table}

\subsection*{Model-theoretic reading}

\begin{enumerate}
\item \textbf{The toggles are independent and the advice is modular.}
Setting several knobs at once is realised by concatenating or overlaying
the corresponding sub-advices --- layout, form factorisation, centre
chain, carry blocks, ring digits are separate stretches. Combining
constructions is thus a Feferman--Vaught-style operation on advices, and
one first-order meta-theorem (linear-time model checking, decidable uniform
theory, and by the advice-automata machinery even FO with cardinality and
Ramsey quantifiers) covers \emph{every} combination at once. Constructions
studied classically in isolation --- extraspecial, Heisenberg, graph
groups, central products --- become a single decidable family.

\item \textbf{The class is closed under its own composition operations},
and those operations are exactly the ``turn on / off'' moves: central
product, direct product, direct factor with an abelian group, base
extension of the centre's exponent. The operations preserving uniform
automaticity \emph{are} the group-theoretic constructions, so the closure
is structural rather than an artefact of the encoding.

\item \textbf{The boundary is sharp and meaningful.} The excluded rows are
precisely the model-theoretically wild objects: Mekler's construction ---
the canonical interpretation of all graphs into $2$-nilpotent groups ---
needs a private target per relation, unbounded-width over every layout; the
ultraspecial extreme is a field-multiplication tensor of full rank-width;
and the relatively free group is Nies--Stephan-non-automatic. Bounded
rank-width is thus the dividing line between the tame class and the
Mekler-wild regime: \emph{our class is the tame complement of Mekler's
construction inside class-$2$ exponent-$p^{d}$ groups.}
\end{enumerate}

\end{document}
